\documentclass{report}
\usepackage{amsmath,amssymb}
\setlength{\parindent}{0mm}
\setlength{\parskip}{1em}
\begin{document}
\begin{center}
\rule{6in}{1pt} \
{\large Steven Hamilton \\
{\bf GPU Implementation of a Monte Carlo Linear Solver Algorithm}}

P O Box 2008 \\ Bldg 5700 \\ Rm H327
\\
{\tt hamiltonsp@ornl.gov}\\
Thomas Evans\\
Stuart Slattery\\
Massimiliano Lupo Pasini\\
Michele Benzi\end{center}

The use of Monte Carlo methods for solving linear systems of equations
has garnered attention in recent years due to the natural resiliency and
scalability of Monte Carlo methods. With most of the largest
supercomputers turning towards vectorized architectures such as GPUs or
the Intel Xeon Phi, we must ask whether Monte Carlo methods are capable
of efficiently utilizing these architectures. At face value, the
stochastic nature of the Monte Carlo approach would appear to lead to
limited vectorization and poor memory access patterns. In this talk, we
discuss an implementation of a Monte Carlo linear solver algorithm on
Nvidia GPUs using the CUDA programming language. We show that careful
organization of data in memory, combined with appropriate use of
GPU-specific capabilities (such as texture fetching), can have a
significant impact on performance.


\end{document}
